% ============================================
% DOC-21
% Extraído de main.pandoc.tex
% ============================================
\fragtitle{doc}{D}{DOC}{21}{Carta de Heitor Oliveira (não enviada)}
\markboth{DOC-21 --- Carta de Heitor Oliveira (não enviada)}{}

\noindent \textit{Documento extraído do arquivo pericial.}

\vspace{0.5em}\noindent\rule{\textwidth}{0.5pt}\vspace{0.5em}

\noindent {\footnotesize\texttt{ARQUIVO P\'UBLICO $\cdot$ PROVA DOCUMENTAL $\cdot$ PROCESSO 1.260}}\\[-0.2em]

\noindent \textbf{Destinatário:} Maria Oliveira (esposa).\\
\noindent \textbf{Remetente:} Dr. Heitor Oliveira.\\
\noindent \textbf{Local:} Barbacena.\\
\noindent\textbf{Data:} não registrada (entre 16 de junho e 20 de junho de 1979).\\
\noindent\textbf{Estado do original:} encontrada dentro da contracapa do caderno pessoal do Dr. Oliveira, dobrada em quatro, sem envelope. Não foi enviada.

\noindent Maria,

\noindent Se você está lendo isto, eu já saí de casa. Ou saí de mim. Não sei qual
\noindent aconteceu primeiro.

\noindent O diário do 1.260 está comigo. Eu não pedi permissão para retirá-lo
\noindent do IML. Assinei o livro de registro como se tivesse autorização, e ninguém
\noindent questionou. É assim que funciona aqui: médico assina, ninguém pergunta.

\noindent Eu achei que ia resolver em dois dias. Pegar o caderno, ler com calma,
\noindent escrever um parecer, devolver. Coisa de rotina. O homem estava morto havia
\noindent poucos dias. O que mais poderia acontecer?

\noindent Maria, eu não li o diário. O diário me leu.

\noindent No terceiro dia, eu acordei com uma mancha amarela no olho. Não dói.
\noindent Não coça. Mas ela está crescendo. O oftalmologista disse que é uma
\noindent hemorragia subconjuntival. Vai sumir em duas semanas. Mas eu sei que não é
\noindent sangue. É a cor do olho que o paciente descreve na página 40. Eu reli a
\noindent página 40 hoje de manhã. A descrição é idêntica à mancha no meu olho.

\noindent No quinto dia, eu senti o cheiro.

\noindent Você conhece o cheiro do Colônia. Você veio comigo uma vez, em 1975.
\noindent Disse que nunca mais voltava. Lembra? O cheiro de creolina, de fezes, de
\noindent comida requentada, de corpo fechado. O diário tem esse cheiro. Mas não veio
\noindent do Colônia. Veio do papel. Como se a tinta tivesse sido feita com alguma
\noindent coisa que apodreceu e foi esquecida.

\noindent No sétimo dia, eu comecei a escrever.

\noindent Não era eu escrevendo. Era a mão. Ela se movia sozinha, igual o
\noindent paciente descreve. Eu via a letra sair --- e não era a minha. Era a letra
\noindent dele. A mesma inclinação para a esquerda. O mesmo \textquotedblleft{}a\textquotedblright{} aberto. A mesma
\noindent pressão irregular no papel, como se a caneta fosse pesada demais.

\noindent Hoje eu acordei e não sabia onde estava. Acordei dentro de uma
\noindent lembrança que não era minha. Uma cerca de arame farpado. O sol queimando
\noindent a nuca. Gente magra demais encostada no arame, com os olhos parados. Eu
\noindent nunca vi aquilo. Mas o corpo lembrava. As mãos tremiam como se tivessem
\noindent segurando arame quente.

\noindent Maria, eu tô com medo de dormir. Toda vez que fecho os olhos, eu
\noindent volto pra vala. É a mesma vala que ele descreve. Os mesmos braços saindo
\noindent da terra. O mesmo cheiro doce que não sai do nariz.

\noindent Eu não vou entregar o diário de volta.

\noindent Eu sei que parece loucura. Mas se eu devolver, alguém vai ler. E se
\noindent alguém ler, a mesma coisa vai acontecer. Eu não posso fazer isso com
\noindent ninguém.

\noindent Ou --- eu não sei se estou protegendo os outros ou se o diário não
\noindent quer ser devolvido.

\noindent Ele gosta de mim.

\noindent Ele gosta de estar comigo.

\noindent Ele não quer voltar pra gaveta.

\noindent Desculpa.

\noindent Heitor

\noindent Nota do arquivista: a carta foi encontrada em 2026, dentro do caderno pessoal do Dr. Oliveira, devolvido ao Arquivo Público Mineiro junto com o diário do paciente 1.260, 47 anos depois de ter sido escrita. A dobra do papel sugere que ele a carregou consigo por todo esse período.

\vspace{0.5em}
\noindent\textit{\textbf{↪ Rotas:} \rota{DOC-01} \rotasep \rota{LUC-02} \rotasep \rota{DOC-22}}